\documentclass{article}
\usepackage[utf8]{inputenc}
\usepackage{listings}
\usepackage{xcolor}
\usepackage{geometry}

% Page layout
\geometry{a4paper, margin=1in}

% Code block styling
\definecolor{codegray}{rgb}{0.95,0.95,0.95}
\definecolor{codeblue}{rgb}{0.0,0.0,0.6}

\lstset{
    backgroundcolor=\color{codegray},
    basicstyle=\ttfamily\small,
    breakatwhitespace=false,
    breaklines=true,
    captionpos=b,
    keepspaces=true,
    showspaces=false,
    showstringspaces=false,
    showtabs=false,
    tabsize=2,
    frame=single,
    keywordstyle=\color{codeblue}\bfseries,
    commentstyle=\color{gray}
}

\title{croSSRoad Database Logic (Pseudocode)}
\author{}
\date{}

\begin{document}

\maketitle

\section*{CrossRoad Database Logic (Pseudocode)}

This document provides a high-level overview of the database schema and access logic for the \textbf{croSSRoad} platform.

\section{Data Model (Schema Structure)}

The database is built on \textbf{PostgreSQL} (via Supabase) and is organized into two main categories: \textbf{Project Management} and \textbf{Scientific Data}.

\subsection{Project Management Tables}
These tables track user submissions and analysis metadata.

\begin{lstlisting}[caption=Project Management Schema]
Table: projects
  - project_id (PK): Unique identifier for the project
  - project_tag:     Unique slug/tag for the project
  - name:            Display name of the project
  - organism:        Organism being analyzed
  - description:     User-provided description
  - created_at:      Timestamp

Table: analysis_runs
  - run_id (PK):     Unique identifier for a specific analysis run
  - project_id (FK): Links to the 'projects' table
  - metrics:         Summary stats (genomes_count, ssrs_count, hotspots_count, etc.)
  - timestamps:      Start/End times for the run
\end{lstlisting}

\subsection{Scientific Data Tables}
These tables store the actual output of the bioinformatics pipeline.

\begin{lstlisting}[caption=Scientific Data Schema]
Table: merged_out (Results from mergedOut.tsv)
  - Stores individual SSR (Simple Sequence Repeat) occurrences per genome.
  - Columns: genome_id, motif, repeat_count, length, location, country, year

Table: hssr_data (Results from hssr_data.csv)
  - Stores Hierarchical SSRs (compound/complex repeats).
  - Columns: genome_id, motif, gene_location, intersecting_gene

Table: mutational_hotspots (Results from mutational_hotspot.csv)
  - Stores regions with high mutation frequency.
  - Columns: motif, gene, variations_list, genome_frequency
\end{lstlisting}

\section{Application Logic (Data Retrieval)}

The application (\texttt{src/routes/croSSRoadDB}) interacts with this database using \textbf{TanStack Query} for state management and the \textbf{Supabase Client} for direct SQL querying.

\subsection{Dashboard Overview (\texttt{index.tsx})}
When a user visits the main database page, the app fetches global statistics in parallel.

\begin{lstlisting}[language=Java, caption=GetDatabaseStats Pseudocode]
FUNCTION GetDatabaseStats():
  // Execute parallel queries for performance
  PARALLEL_START
    projects_count    = COUNT(*) FROM projects
    runs_count        = COUNT(*) FROM analysis_runs
    ssrs_count        = COUNT(*) FROM merged_out
    hssrs_count       = COUNT(*) FROM hssr_data
    hotspots_count    = COUNT(*) FROM mutational_hotspots
    
    // Fetch recent activity
    recent_projects   = SELECT * FROM projects ORDER BY created_at DESC LIMIT 5
    recent_runs       = SELECT * FROM analysis_runs ORDER BY timestamp DESC LIMIT 5
    
    // Fetch sample data for distinct counts (Country/Year/Motif stats)
    stats_sample      = SELECT * FROM merged_out LIMIT 1000
  PARALLEL_END

  RETURN {
    total_projects: projects_count,
    total_genomes:  SUM(genomes_used_count) FROM analysis_runs,
    total_ssrs:     ssrs_count,
    unique_countries: CALCULATE_UNIQUE(stats_sample.country),
    unique_years:     CALCULATE_UNIQUE(stats_sample.year)
  }
END FUNCTION
\end{lstlisting}

\subsection{Search Functionality (\texttt{search.tsx})}
The search bar performs a ``federated search'' across multiple tables simultaneously.

\begin{lstlisting}[language=Java, caption=SearchDatabase Pseudocode]
FUNCTION SearchDatabase(user_query):
  // 1. Search Projects
  results_projects = SELECT * FROM projects 
                     WHERE name LIKE %user_query% 
                     OR organism LIKE %user_query%
                     OR description LIKE %user_query%

  // 2. Search SSRs
  results_ssrs = SELECT * FROM merged_out 
                 WHERE motif LIKE %user_query% 
                 OR genome_id LIKE %user_query%
                 OR country LIKE %user_query%
                 LIMIT 50

  // 3. Search Hotspots
  results_hotspots = SELECT * FROM mutational_hotspots 
                     WHERE motif LIKE %user_query% 
                     OR gene LIKE %user_query%
                     LIMIT 50

  RETURN COMBINE(results_projects, results_ssrs, results_hotspots)
END FUNCTION
\end{lstlisting}

\end{document}
